\subsection{庞加莱群胚的无局部单值作用}

引理 1. --- 设 B 是局部道路连通的拓扑空间，\(p \colon E \to B\) 和 \(p' \colon E' \to B\) 是满足道路提升性质的平展分离映射（III，第 302 页）。设 \(g \colon E \to E'\) 是满足 \(p' \circ g = p\) 的映射。假设 \(g\) 与群胚 \(\varpi(B)\) 在 \(E\) 和 \(E'\) 上的典范作用相容。则 \(g\) 连续。

设 c 是 E 中的一条道路；置 \(c' = g \circ c\)，并证明映射 \(c'\) 连续。记 B 中的道路 \(p \circ c\) 为 d，并对每个 \(t \in I\)，记 \(d_t\) 为道路 \(s \mapsto d(st)\)。对于所有 \(t \in I\)，有 \(c(t) = c(0) \cdot d_t\)（III，第 304 页，注记 3），由关于 g 的假设，从而有 \(c'(t) = c'(0) \cdot d_t\)，这根据同一注记说明 \(c'\) 是道路 d 的提升。这说明映射 g 按道路连续。由于空间 E 局部道路连通（III，第 261 页，命题 2 的推论 2），映射 g 连续（III，第 269 页，命题 13 的推论）。

设 B 是拓扑空间。考虑一种作用 \(\varphi = (\varphi_{a,b})_{(a,b)\in \mathrm{B}\times \mathrm{B}}\)，它是群胚 \(\varpi (\mathrm{B})\) 关于映射 \(p\colon \mathrm{E}\to \mathrm{B}\) 在集合 E 上的作用。称 \(\varpi (\mathrm{B})\) 在 E 上无单值性地作用，若对每个 \(b\in \mathrm{B}\) 和每个圈类 \(\gamma \in \pi_1(\mathrm{B},b)\)，\(\gamma\) 在纤维 \(\mathrm{E}_b\) 上的作用是平凡的，则称作用 \(\varphi\)（即群胚 \(\varpi (\mathrm{B})\) 的作用）无单值性地作用在 E 上（参见 II，第 168 页）。若 B 道路连通，只需对 B 的一个点验证这一点即可（同上）。如果 B 的每一点都有邻域 V，使得 \(\varpi (\mathrm{V})\) 在集合 \(\mathrm{E_V} = p^{-1}(\mathrm{V})\) 上关于映射 \(p_{\mathrm{V}} = p\mid p^{-1}(\mathrm{V})\) 无单值性地作用，则称群胚的作用无局部单值性。

注记。--- 设 B 是局部道路连通的拓扑空间，并假设 B 的每一点 b 都有邻域 V，使得 \(\pi_{1}(V,b)\) 在 \(\pi_{1}(B,b)\) 中的像约化为单位元（参见 IV，第 340 页，定义 2）。那么，群胚 \(\varpi(B)\) 的每个作用都无局部单值性。

事实上，考虑集合 E、映射 \(p \colon \mathrm{E} \to \mathrm{B}\) 以及作用 \(\varphi\)；它是群胚 \(\varpi(\mathrm{B})\) 关于 \(p\) 在 E 上的作用。设 \(a\) 是 B 的一点，V 是 \(a\) 的邻域，并且 \(\pi_1(\mathrm{V}, a)\) 在 \(\pi_1(\mathrm{B}, a)\) 中的像约化为单位元。设 U 是包含于 V 的 \(a\) 的道路连通邻域。设 \(b\) 是 U 的一点，\(\gamma \in \pi_1(\mathrm{U}, b)\)。设 \(\delta\) 是 U 中连接 \(a\) 与 \(b\) 的道路的类。于是，\(\delta\gamma\delta^{-1}\) 是 U 中以 \(a\) 为基点的圈的类；因此，它在 \(\pi_1(\mathrm{B}, a)\) 中的像是平凡的。从而有 \(\varphi_{a,a}(\delta\gamma\delta^{-1}) = \mathrm{Id}_{\mathrm{E}_a}\)，所以 \(\varphi_{b,b}(\gamma) = \mathrm{Id}_{\mathrm{E}_b}\)。这说明 \(\varpi(\mathrm{U})\) 关于 U-空间 \(\mathrm{E_U}\) 的作用无单值性。因此，作用 \(\varphi\) 无局部单值性。

命题 3. --- 设 B 是局部道路连通的拓扑空间，E 是赋有作用 \(\varphi\)（群胚 \(\varpi(B)\) 的作用）的集合，该作用关于映射 \(p: E \to B\) 无局部单值性。于是 E 上存在唯一的拓扑，使得下列条件成立：

\begin{enumerate}
\def\labelenumi{(\roman{enumi})}
\item
  赋予此拓扑的集合 E 连同映射 p 是 B 的覆叠；
\item
  \(\varpi(\mathrm{B})\) 在此覆叠上的典范作用与作用 \(\varphi\) 相同。
\end{enumerate}

这种拓扑的唯一性由 III，第 312 页的引理 1 得出，其中取 E 的恒等映射作为 g。为证明其存在性，还可假设 B 连通且非空。

首先假设作用 \(\varphi = (\varphi_{a,b})_{(a,b)\in \mathrm{B}\times \mathrm{B}}\) 是群胚 \(\varpi (\mathrm{B})\) 关于 \(p\) 在集合 E 上的无单值作用。

设 \(a\) 是 B 的一点。对于每个 \(b \in \mathrm{B}\)，存在一条道路 \(c\) 在 B 中连接点 \(a\) 与点 \(b\)。若 \(\gamma \in \varpi_{a,b}(\mathrm{B})\) 表示道路 \(c\) 的类，则双射 \(\varphi_{a,b}(\gamma) \colon \mathrm{E}_a \to \mathrm{E}_b\) 与道路 \(c\) 无关，因为该作用无单值性；将此双射记为 \(f_{a,b}\)。映射 \(\Phi_a \colon \mathrm{B} \times \mathrm{E}_a \to \mathrm{E}\)（由 \((b,x) \mapsto f_{a,b}(x)\) 定义）是双射；其逆双射将 \(x \in \mathrm{E}\) 映到二元组 \((p(x), f_{p(x),a}(x))\)。赋予集合 \(\mathrm{E}_a\) 离散拓扑，于是 B-空间 \(\mathrm{B} \times \mathrm{E}_a\) 是 B 的平凡覆叠。通过结构转移，双射 \(\Phi_a\) 在 E 上赋予一个拓扑，使其成为 B 的覆叠。

证明 \(\varpi(\mathrm{B})\) 在 E 上的典范作用与作用 \(\varphi\) 相同。设 \(x\) 是 \(\mathrm{E}_a\) 的一点，设 \(b\) 是 B 的一点。设 \(c\) 是 B 中连接点 \(a\) 与点 \(b\) 的道路；映射 \(t \mapsto \Phi_a(c(t), x)\) 是 E 中连接点 \(x\) 与点 \(\Phi_a(b, x) = f_{a,b}(x)\) 的道路。若 \(\gamma \in \varpi_{a,b}(\mathrm{B})\) 表示道路 \(c\) 的类，则有 \(x \cdot \gamma = f_{a,b}(x)\)。这说明 \(\varpi(\mathrm{B})\) 在覆叠 E 上的典范作用与作用 \(\varphi\) 在起点为 \(a\) 的道路类上相同。由于这些类生成群胚 \(\varpi(\mathrm{B})\)，两个作用相同。

现在处理一般情形。令 \(\mathcal{B}\) 为 E 的开子集 V 的集合，其中 \(\varpi(\mathrm{V})\) 在 \(p^{-1}(\mathrm{V})\) 上无单值性地作用。根据假设，\(\mathcal{B}\) 中的元素覆盖 B。由上文，对于每个 \(\mathrm{V} \in \mathcal{B}\)，在集合 \(p^{-1}(\mathrm{V})\) 上存在唯一拓扑，使得 \((p^{-1}(\mathrm{V}), p_{\mathrm{V}})\) 成为 V 的覆叠，并且 \(\varpi(\mathrm{V})\) 在此覆叠上的典范作用与 \(\varphi\) 诱导的作用相同。

设 \(V' \in \mathcal{B}\)。上面定义的在 \(p^{-1}(V \cap V')\) 上由 \(p^{-1}(V)\) 的拓扑（分别由 \(p^{-1}(V')\) 的拓扑）诱导的拓扑，确实使其成为 \(V \cap V'\) 的覆叠，并且 \(\varpi(V \cap V')\) 在其上的典范作用由作用 \(\varphi\) 诱导。因此这些拓扑都与 \(p^{-1}(V \cap V')\) 的拓扑相同。于是存在 E 上唯一的拓扑，在每个 \(p^{-1}(V)\) 上诱导出先前定义的拓扑（参见 I，§2，第 16 页）。

赋予 E 此拓扑后，映射 p 连续，且 B-空间 E 是覆叠。\(\varpi(\mathrm{B})\) 的典范作用与作用 \(\varphi\) 在像包含于 \(\mathcal{B}\) 中某个开集内的道路类上相同。由引理 4 的

III，第 272 页，这些类生成群胚 \(\varpi(B)\)。因此这两个作用相同（II，第 167 页）。

推论。--- 设 B 是连通且局部道路连通的拓扑空间，(E,p) 是 B-空间，其投影 p 是平展、分离且具有道路提升性质的映射。若 \(\varpi(\mathrm{B})\) 关于 p 在集合 E 上的典范作用无局部单值性，则 E 是 B 的覆叠。

给 E 赋予由道路提升定义的 \(\varpi(B)\) 的典范作用（III，第 303 页，编号 3）。E 上存在一个拓扑，使条件 (i) 和 (ii) 的命题 3 得以满足。根据 III，第 312 页的引理 1，此拓扑与 E 上给定的拓扑相同。

